\chapter[Band 20: Endliche Mengen]{Band 20 auf einen Blick}
\noindent\textbf{Endliche Mengen}\par\medskip
\noindent\emph{Lesart.}
\(\Finite M\) ist das primitive Endlichkeitsprädikat; seine
Bijektionscharakterisierung ist ein Satz.
\(\NatSeg n=\{k\in\mathbb N\mid k<n\}\) bezeichnet einen
Peano-Anfangsabschnitt und \(\EqCard\) die Gleichmächtigkeit durch Bijektion.
Die im Original ausdrücklich angegebenen Auswahlabhängigkeiten der
Dedekind-/Tarski- und Unendlichkeitscharakterisierungen gelten auch hier.

\begin{BandUebersicht}
\textbf{Axiome: leere Menge und Adjunktion}
\UebFormel{\begin{aligned}
&\Finite\varnothing,\\
&\Finite A,\ a\notin A\vdash\Finite{A\cup\{a\}}.
\end{aligned}}
\UebQuellen{\FormulaRefAuto{Endlichkeitsaxiome}[def];
\FormulaRefAuto{FiniteEmptyAxiom}[ax];
\FormulaRefAuto{FiniteAdjunctionAxiom}[ax]}
&
\textit{Abschlüsse der Endlichkeit}
\UebFormel{\begin{aligned}
&\Finite A\vdash\Finite{A\cup\{a\}},\\
&\Finite A,\ \Finite B\vdash\Finite{A\cup B},\\
&B\subseteq A,\ \Finite A\vdash\Finite B,\\
&F\colon A\to B,\ \Finite A\vdash\Finite{F[A]}.
\end{aligned}}
\UebQuellen{\FormulaRefAuto{FiniteInsert}[thm];
\FormulaRefAuto{FiniteUnion}[thm];
\FormulaRefAuto{FiniteSubset}[thm];
\FormulaRefAuto{FiniteImage}[thm]}
\\ \addlinespace[.8ex]

\textbf{Axiom: Induktion über endliche Mengen}
Aus dem Anfang \(P(\varnothing)\) und dem Schritt
\UebFormel{\begin{aligned}
&\Finite A,\ P(A),\ a\notin A\\
&\qquad\vdash P(A\cup\{a\})
\end{aligned}}
folgt unter der Voraussetzung \(\Finite M\) die Aussage \(P(M)\).
\UebQuellen{\FormulaRefAuto{FiniteInductionRule}[ax]}
&
\textit{Bijektionscharakterisierung}
\UebFormel{\begin{aligned}
&n\in\mathbb N\vdash\Finite{\NatSeg n},\\
&\Finite M\eqvdash
\exists n\in\mathbb N\;(\NatSeg n\EqCard M).
\end{aligned}}
Für die von-Neumann-Realisierung gilt außerdem
\UebFormel{\Finite M\eqvdash
\exists n\in\mathbb N\;(n\EqCard M).}
\UebQuellen{\FormulaRefAuto{PeanoNatSegFinite}[thm];
\FormulaRefAuto{FiniteEqCardEquiv}[thm];
\FormulaRefAuto{FiniteVonNeumannEqCardEquiv}[thm]}
\\ \addlinespace[.8ex]

\textbf{Aufzählungsmengen}
Für \(a\colon\mathbb N\to A\) und \(n\in\mathbb N\):
\UebFormel{\operatorname{Enum}_a(n)\coloneqq a[\NatSeg n].}
\UebQuellen{\FormulaRefAuto{EnumSetDef}[def]}
&
\textit{Rekursion und Elementkriterium}
\UebFormel{\begin{aligned}
&\operatorname{Enum}_a(0)=\varnothing,\\
&\operatorname{Enum}_a(n+1)
=\operatorname{Enum}_a(n)\cup\{a_n\},\\
&x\in\operatorname{Enum}_a(n)\\
&\qquad\leftrightarrow\exists i\in\NatSeg n\;(x=a_i).
\end{aligned}}
\UebQuellen{\FormulaRefAuto{EnumSetZero}[thm];
\FormulaRefAuto{EnumSetSucc}[thm];
\FormulaRefAuto{EnumSetMembership}[thm]}
\\ \addlinespace[.8ex]

\textbf{Tarski-Endlichkeit}
\(\TarskiFinite(M)\) fordert für jede
\(\varnothing\neq B\subseteq\powerset(M)\):
\UebFormel{\exists x\in B\;\forall y\in B\;
(x\subseteq y\to y=x).}
\(\TarskiMinFinite(M)\) ersetzt hier \(x\subseteq y\) durch
\(y\subseteq x\).
\UebQuellen{\FormulaRefAuto{Tarski-endliche Menge}[def];
\FormulaRefAuto{minimal Tarski-endliche Menge}[def]}
&
\textit{Äquivalente Endlichkeitsbegriffe}
\UebFormel{\begin{aligned}
&\Finite M\leftrightarrow
\neg\exists H\colon\mathbb N\inj M,\\
&\Finite M\leftrightarrow\TarskiFinite(M),\\
&\TarskiFinite(M)\leftrightarrow\TarskiMinFinite(M).
\end{aligned}}
Die nichttriviale Rückrichtung aus dem Injektionsausschluss
verwendet im Band das Auswahlaxiom.
\UebQuellen{\FormulaRefAuto{\Finite{M}\eqvdash
\neg\exists H\colon\mathbb{N}\inj M}[thm];
\FormulaRefAuto{\Finite{M}\eqvdash\TarskiFinite(M)}[thm];
\FormulaRefAuto{TarskiMaxMinFiniteEquiv}[thm]}
\\ \addlinespace[.8ex]

\textbf{Kardinalzahl einer endlichen Menge}
Unter \(\Finite M\):
\UebFormel{\begin{aligned}
\FiniteCard(M)\coloneqq\iota n\;&
(n\in\mathbb N\\
&{}\land\NatSeg n\EqCard M).
\end{aligned}}
\UebQuellen{\FormulaRefAuto{FiniteCardinality}[def]}
&
\textit{Eindeutigkeit und Zählung}
\UebFormel{\begin{aligned}
&\exists!n\in\mathbb N\;(\NatSeg n\EqCard M),\\
&\FiniteCard(M)\in\mathbb N,\\
&\NatSeg{\FiniteCard(M)}\EqCard M.
\end{aligned}}
Die Definition wählt somit einen nachweislich eindeutigen
natürlichen Zählwert.
\UebQuellen{\FormulaRefAuto{FiniteCardWitnessUnique}[thm];
\FormulaRefAuto{FiniteCardWellDefined}[thm];
\FormulaRefAuto{FiniteCardNatural}[thm];
\FormulaRefAuto{FiniteCardCounts}[thm]}
\\ \addlinespace[.8ex]

\textbf{Binäre Teilmengencodierung}
Für \(X\subseteq\NatSeg n\):
\UebFormel{\begin{aligned}
\varepsilon_X(i)&\coloneqq
\begin{cases}1,&i\in X,\\0,&i\notin X,\end{cases}\\
c_n(X)&\coloneqq\sum_{i<n}\varepsilon_X(i)2^i.
\end{aligned}}
\UebQuellen{\FormulaRefAuto{BinarySubsetCode}[def]}
&
\textit{Potenzmengen zählen}
\UebFormel{c_n\colon\powerset(\NatSeg n)\bij\NatSeg{2^n}.}
Für \(\Finite A\) und \(\FiniteCard(A)=n\):
\UebFormel{\begin{aligned}
&\FiniteCard(\powerset(A))=2^n,\\
&\FiniteCard(\PowNE A)=2^n-1.
\end{aligned}}
Dabei ist \(\PowNE A\) die Menge der nichtleeren Teilmengen.
\UebQuellen{\FormulaRefAuto{BinarySubsetCodeBijective}[thm];
\FormulaRefAuto{FinitePowerSetCardinality}[thm];
\FormulaRefAuto{FiniteNonemptyPowerSetCardinality}[thm]}
\\ \addlinespace[.8ex]

\textbf{Maximalelementeigenschaft}
\UebFormel{\begin{aligned}
&\operatorname{HatMax}_{A,\preceq}(C)\coloneqq\\
&\quad(C\subseteq A\land C\neq\varnothing\\
&\qquad\to\exists m\in C\;\forall b\in C\\
&\qquad\qquad(m\preceq b\to b=m)).
\end{aligned}}
\UebQuellen{\FormulaRefAuto{FinitePosetMaximalProperty}[def]}
&
\textit{Endliche partielle Ordnungen}
Unter \(\PartOrd{A,\preceq}\),
\(\varnothing\neq B\subseteq A\) und \(\Finite B\):
\UebFormel{\begin{aligned}
&\exists m\in B\;\forall b\in B\;(m\preceq b\to b=m),\\
&\exists \ell\in B\;\forall b\in B\;(b\preceq\ell\to b=\ell).
\end{aligned}}
Ein maximales bzw. minimales Element muss kein größtes bzw.
kleinstes Element sein.
\UebQuellen{\FormulaRefAuto{FinitePosetMaximalElement}[thm];
\FormulaRefAuto{FinitePosetMinimalElement}[thm]}
\\ \addlinespace[.8ex]

\textbf{Unendlich, abzählbar und überabzählbar}
\UebFormel{\begin{aligned}
\Infinite A&\coloneqq\neg\Finite A,\\
\AtMostCountable A&\coloneqq A\CardLeq\mathbb N,\\
\CountablyInfinite A&\coloneqq A\EqCard\mathbb N,\\
\Uncountable A&\coloneqq\neg\AtMostCountable A.
\end{aligned}}
\UebQuellen{\FormulaRefAuto{InfiniteDef}[def];
\FormulaRefAuto{AtMostCountableDef}[def];
\FormulaRefAuto{CountablyInfiniteDef}[def];
\FormulaRefAuto{UncountableDef}[def]}
&
\textit{Unendlichkeit und abzählbare Alternative}
\UebFormel{\begin{aligned}
&\Infinite{\mathbb N},\\
&\Infinite A\leftrightarrow\exists H\colon\mathbb N\inj A,\\
&\AtMostCountable A\leftrightarrow\\
&\qquad(\Finite A\lor\CountablyInfinite A).
\end{aligned}}
Für die angegebenen Rückführungen auf Injektionen wird die im Band
festgehaltene Auswahlvoraussetzung benutzt.
\UebQuellen{\FormulaRefAuto{NaturalNumbersInfinite}[thm];
\FormulaRefAuto{InfiniteInjectionCharacterization}[thm];
\FormulaRefAuto{AtMostCountableDichotomy}[thm]}
\\ \addlinespace[.8ex]

\textbf{Einermengenabbildung und Diagonalmenge}
\UebFormel{\begin{aligned}
&\SingletonEmbedding A(x)\coloneqq\{x\},\\
&F\colon A\to\powerset(A)\vdash\\
&\qquad D_F^A\coloneqq\{x\in A\mid x\notin F(x)\}.
\end{aligned}}
\UebQuellen{\FormulaRefAuto{SingletonEmbeddingDef}[def];
\FormulaRefAuto{CantorDiagonalSetDef}[def]}
&
\textit{Cantors strikter Potenzmengensatz}
\UebFormel{\begin{aligned}
&\SingletonEmbedding A\colon A\inj\powerset(A),\\
&x\in A\vdash F(x)\neq D_F^A,\\
&A\CardLt\powerset(A).
\end{aligned}}
Insbesondere gibt es keine größte Mächtigkeitsstufe.
\UebQuellen{\FormulaRefAuto{SingletonPowerSetInjection}[thm];
\FormulaRefAuto{CantorDiagonalMissed}[thm];
\FormulaRefAuto{CantorStrictPowerSet}[thm];
\FormulaRefAuto{NoGreatestCardinality}[thm]}
\\ \addlinespace[.8ex]

\textbf{Kontinuumshypothese als Aussage}
\UebFormel{\begin{aligned}
\mathsf{CH}\coloneqq\neg\exists A\;(&\mathbb N\CardLt A\\
&{}\land A\CardLt\powerset(\mathbb N)).
\end{aligned}}
Dies ist die Definition einer Hypothese; der Band beweist sie nicht.
\UebQuellen{\FormulaRefAuto{ContinuumHypothesisRelational}[def]}
&
\textit{Bewiesene Unendlichkeitsstufen}
\UebFormel{\begin{aligned}
&\mathbb N\CardLt\powerset(\mathbb N),\\
&\powerset(\mathbb N)\CardLt
\powerset(\powerset(\mathbb N)),\\
&\mathbb N\CardLt\mathbb R,\qquad\Uncountable{\mathbb R}.
\end{aligned}}
Keine Folge natürlicher Indizes zählt sämtliche reellen Zahlen auf.
\UebQuellen{\FormulaRefAuto{FirstInfinitePowerSetLevels}[thm];
\FormulaRefAuto{NaturalsStrictlyLessPowerfulThanReals}[thm];
\FormulaRefAuto{RealsUncountable}[thm];
\FormulaRefAuto{CantorNoRealSurjection}[thm]}
\\
\end{BandUebersicht}
